 % ============================================================
%  SEKCJA 8: WNIOSKI
% ============================================================

\section{Wnioski}

Niniejszy artykuł podjął się analizy stanowiska
Apimondia dotyczącego produkcji i~dystrybucji miodu
niedojrzałego, ze~szczególnym uwzględnieniem implikacji
mikrobiologicznych, regulacyjnych i~rynkowych stosowania
liczby drożdży osmofilnych oraz kryterium morfologicznego
zasklepienia komórek plastra jako wskaźników jakości.
Przeprowadzona analiza prowadzi do~szeregu wniosków
o~różnym stopniu uogólnienia, od~szczegółowych
konkluzji biochemicznych po~rekomendacje normatywne
o~zasięgu globalnym.

%,----------------------------------------------------------
\subsection{Wnioski naukowe}
%,----------------------------------------------------------

\textbf{(W1) Obecność drożdży osmofilnych w~miodzie jest zjawiskiem biologicznie normalnym.}
Drożdże z~rodzajów \textit{Zygosaccharomyces}, \textit{Schizosaccharomyces}
i~\textit{Metschnikowia} są nieodłącznym składnikiem ekosystemu
ula, niezależnie od~metody produkcji, systemu ula i~kraju
pochodzenia. Naturalna liczebność populacji drożdży
w~prawidłowym, niefermentującym miodzie komercyjnym
może wynosić kilkaset do~kilku tysięcy CFU/g bez~jakiegokolwiek
wpływu na~stabilność mikrobiologiczną produktu, o~ile
aktywność wody $a_w < 0{,}60$~\cite{ruegg1981}.
Wniosek ten jest ugruntowany w~literaturze od~dekad
i~potwierdzony przez wszystkie aktualne przeglądy
systematyczne~\cite{biochemreactions2022,rodriguezmachado2024}.

\textbf{(W2) Aktywność wody, nie wilgotność ani liczba drożdży, jest pierwotnym parametrem ryzyka fermentacji.}
Fizjologiczny próg wzrostu drożdży osmofilnych
($a_w = 0{,}61$--$0{,}62$) wyznacza biologicznie
uzasadnioną granicę, poniżej której fermentacja miodu
jest termodynamicznie niemożliwa, niezależnie od~liczby
drożdży wyrażonej w~CFU/g~\cite{waaw2006,ruegg1981}.
Krystalizacja glukozy stanowi dodatkowy czynnik dynamicznie
modyfikujący $a_w$ fazy ciekłej, co~sprawia, że~ocena
ryzyka fermentacji musi uwzględniać stan fazowy produktu
i~nie~może być oparta wyłącznie na~globalnej zawartości
wody~\cite{zamora2006,ijfans2019}.

\textbf{(W3) Dojrzewanie miodu jest ciągłym procesem biologicznym, a~nie zdarzeniem binarnym.}
Wskaźniki dojrzałości, aktywność inwertazy, oksydazy
glukozowej, diastazy, profil kwasów organicznych,
stosunek fruktozy do~glukozy, zmieniają się gradientowo
w~trakcie procesu dojrzewania i~nie korelują jednoznacznie
z~momentem zasklepienia komórek przez pszczoły~\cite{zhang2021,sun2021}.
Miód niezasklepiony może wykazywać profil chemiczny
właściwy miodowi dojrzałemu, jeśli wilgotność jest niska;
miód zasklepiony może utracić cechy dojrzałości wskutek
nieodpowiedniej obróbki termicznej.~\cite{codex_stan12,guo2020,zhang2021}.

\textbf{(W4) Standardy miodowe mają charakter eurocentryczny i~są nieadekwatne dla globalnych systemów produkcji oraz uwarunkowań mikroklimatycznych.}
Klimatyczne, technologiczne i~biologiczne uwarunkowania
produkcji miodu w~strefach tropikalnych i~subtropikalnych
, w~tym wysoka wilgotność środowiskowa, systemy uli
kłodowych i~koszowych, odmienne właściwości miodów
pszczół bezżądłowych (Meliponini), sprawiają,
że~jednolite kryterium morfologiczne (zasklepienie)
i~mikrobiologiczne (CFU/g bez~$a_w$) nie~mogą pełnić
roli globalnych wskaźników jakości bez~systematycznego
faworyzowania miodów produkowanych w~systemach
europejskich~\cite{tadele2025,stinglessyeastbrazil2021,sciendo2012}.
Stanowisko to~jest zbieżne z~ogólnym trendem w~prawodawstwie
żywnościowym w~kierunku uznawania różnorodności lokalnych
systemów produkcji i~odejścia od~norm wywodzących się
wyłącznie z~praktyk zachodnioeuropejskich~\cite{phiri2022}.

\textbf{(W5) Dominującą formą fałszowania miodu na~rynku UE jest adulteracja syropami cukrowymi, nie dystrybucja miodu o~podwyższonej liczbie drożdży.}
Dane skoordynowanej akcji kontrolnej Komisji Europejskiej
,,From the Hives'' (2021--2022) jednoznacznie wskazują,
że~46\% podejrzanych przesyłek importowanego miodu
zawierało dodatki syropów C$_3$/C$_4$
lub~oligosacharydów~\cite{eufromhives2021}. Żaden
z~dostępnych raportów kontrolnych nie~identyfikuje
podwyższonej liczby CFU/g jako głównego wskaźnika
fałszowania w~praktyce inspekcyjnej. Alokacja zasobów
kontrolnych w~kierunku pomiaru $a_w$ jest
zatem bardziej adekwatna do~rzeczywistego profilu
ryzyka rynkowego niż intensyfikacja metod
mikrobiologicznych~\cite{schoder2026,eufromhives2021}.

%,----------------------------------------------------------
\subsection{Ocena stanowiska Apimondia}
%,----------------------------------------------------------

Stanowisko Apimondia z~2023~roku w~sprawie miodu niedojrzałego
jest wartościowym i~potrzebnym wkładem w~ochronę integralności
rynku pszczelarskiego. Zagrożenie ze~strony miodu systemowo
niedojrzałego, produkowanego metodą przemysłowego
odwadniania nektaru w~krótkim czasie i~bez~zachowania
biologicznych procesów dojrzewania, jest realne i~stanowi
nieuczciwe konkurowanie z~producentami przestrzegającymi
standardów jakości. W~tym zakresie intencja i~kierunek
działania Apimondia zasługują na~pełne poparcie.

Niniejszy artykuł nie~podważa tej intencji, lecz wskazuje,
że~obecne sformułowanie stanowiska jest analitycznie
niekompletne w~dwóch istotnych obszarach. Po~pierwsze,
nie~rozróżnia ono między (a)~fermentacją jako cechą
produktu w~chwili zbioru a~(b)~fermentacją jako zdarzeniem
post-harvest wynikającym z~błędów w~łańcuchu dostaw.
Po~drugie, nie~uwzględnia fundamentalnej roli aktywności
wody jako parametru nadrzędnego wobec liczby drożdży
w~ocenie ryzyka fermentacyjnego. Uzupełnienie stanowiska
Apimondia o~te dwa elementy nie~osłabi jego przekazu
ochronnego, wzmocni je~natomiast naukowo
i~uczyni bardziej odpornym na~prawne kwestionowanie
w~sporach handlowych z~importerami z~krajów
tropikalnych~\cite{ruegg1981,tadele2025,analytica2020}.

Po trzecie , może powodować dyskryminację miodów ( w tym europejskich), dla której naturalnie wysoka zawartość drożdży osmofilnych jest częsta/charakterystyczna, a której to wartości nie są spowodowane fermentacją. 

%,----------------------------------------------------------
\subsection{Propozycja modelu: synteza}
%,----------------------------------------------------------

Proponowany model oceny dojrzałości funkcjonalnej (MODF)
oferuje operacyjne narzędzie oceny jakości miodu, które:

\begin{itemize}
    \item jest \textbf{naukowo uzasadnione}, oparte
          na~termodynamicznych właściwościach aktywności
          wody, enzymatycznej kinetyce dojrzewania
          i~metabolomice~\cite{ruegg1981,zhang2021,guo2020};
    \item jest \textbf{zgodne z~obowiązującymi regulacjami}
         , Codex STAN 12-1981 i~Dyrektywa UE 2001/110/WE
         , i~nie~wymaga ich rewizji, jedynie interpretacyjnego
          rozszerzenia~\cite{codex_stan12,eudir2001110};
    \item jest \textbf{technicznie dostępne}, pomiar $a_w$
          i~refraktometria są realizowalne w~każdym
          laboratorium kontrolnym bez~specjalistycznej
          infrastruktury~\cite{waaw2006,zamora2006};
    \item jest \textbf{globalnie stosowalny}, kryteria oparte
          na~$a_w$ i~profilu enzymatycznym są niezależne
          od~systemu ula i~warunków klimatycznych
          miejsca produkcji~\cite{tadele2025,stinglessyeastbrazil2021};
    \item jest \textbf{proporcjonalny do~ryzyka},
          zapobiega eskalacji do~kosztownych metod.~\cite{schoder2026,nmrhoney2023}.
\end{itemize}

%,----------------------------------------------------------
\subsection{Ograniczenia i kierunki dalszych badań}
%,----------------------------------------------------------

Autorzy artykułu zdają sobie sprawę z~ograniczeń
niniejszego opracowania. Proponowany model MODF opiera
się na~danych empirycznych dostępnych w~momencie
przygotowania manuskryptu i~wymaga dalszej weryfikacji
w~ramach prospektywnych badań wielolaboratoryjnych.
Wskazane są w~szczególności następujące kierunki
dalszych badań:

\begin{enumerate}
    \item \textbf{Walidacja wielolaboratoryjna biomarkera
          kwasu dekenodwukarboksylowego} na~próbkach
          reprezentujących globalne zróżnicowanie typów
          botanicznych i~systemów produkcji, wykraczając
          poza trzy typy (rzepak, akacja, jujube) zbadane
          przez Sun i~in.~\cite{sun2021}; walidacja
          powinna obejmować miody bezżądłowe (Meliponini)
          oraz miody z~systemów tradycyjnych Afryki
          Wschodniej~\cite{tadele2025};

%     \item \textbf{Budowa globalnych baz referencyjnych
%           NMR i~metabolomicznych} reprezentujących
%           różnorodność geograficzną i~botaniczną produkcji
%           miodu poza Europą; bez~tych baz metody omiczne
%           pozostają narzędziami o~ograniczonej mocy
%           atrybucji dla~miodów z~krajów
%           rozwijających się~\cite{schoder2026,tsagkaris2021};

    \item \textbf{Opracowanie i~walidacja wielowymiarowego
          wskaźnika dojrzałości funkcjonalnej} (FMI,
          \textit{Functional Maturity Index}) integrującego
          wyniki parametrów Poziomu~II w~jeden wskaźnik
          o~wartości 0--100, kalibrowanego dla~różnych
          typów botanicznych i~geograficznych; taki wskaźnik
          mógłby zastąpić binarne decyzje dyskwalifikacyjne
          podejściem gradientowym, proporcjonalnym
          do~odstępstwa od~stanu optymalnego~\cite{zhang2021,guo2020,sun2021};

%     \item \textbf{Prospektywna analiza łańcucha dostaw}
%           dla~miodów importowanych z~krajów tropikalnych,
%           rejestrująca parametry środowiskowe (T, RH)
%           w~całym cyklu od~pasieki do~punktu kontroli
%           granicznej; tylko takie dane pozwolą na~empiryczne
%           rozróżnienie fermentacji produkcyjnej
%           od~post-harvest i~na~walidację postulowanego
%           protokołu traceability~\cite{analytica2020,apinz2021,cbi_traceability2024};

    \item \textbf{Harmonizacja standardów ISO TC~34/SC~17
          z~uwzględnieniem systemu produkcji meliponikultury}
         , kategoria miodów bezżądłowych wymaga odrębnego
          toru normalizacyjnego z~własnymi progami
          wilgotności, $a_w$ i~kryteriami enzymatycznymi,
          uwzględniającymi biologiczne właściwości
          tych owadów~\cite{stinglessyeastbrazil2021,stinglesswjava2024}.
\end{enumerate}

Podjęcie tych badań jest możliwe wyłącznie w~ramach
współpracy międzynarodowej obejmującej instytucje naukowe,
organizacje pszczelarskie i~organy regulacyjne z~krajów
zarówno ,,globalnej Północy'', jak i krajów tropikalnych będących głównymi producentami miodu. Dialog ten jest nie~tylko potrzebą naukową,
lecz warunkiem koniecznym budowy globalnie sprawiedliwego i~naukowo wiarygodnego systemu oceny jakości jednego
z~najstarszych produktów spożywczych wytwarzanych przez człowieka.